\section{Results}
\label{sec:results}

We present the numerical results of our replication study, which confirm the complex phase behavior of the emergent magnetic monopole lattice. The results are obtained by maximizing the free energy functional $G$ across the mixing ratio $p$, magnetic field $h$, and temperature $T$.

\subsection{h-T Phase Diagrams}
For a fixed mixing ratio $p=0.4$, we compute the phase maps in the $h-T$ plane for external fields applied along the $[100]$, $[110]$, and $[111]$ directions. As shown in Fig.~\ref{fig:phase_ht}, the system exhibits a rich variety of phases. At low fields, the 3Q state dominates, while increasing $h$ drives transitions toward the 1Q phase. The phase boundaries are characterized by jumps in the monopole charge $N_m$ and anomalies in the magnetization $m$.

\subsection{p-h Ground State Sequence}
The influence of the mixing ratio $p$ on the ground state is captured in the $p-h$ phase diagram (Fig.~\ref{fig:phase_ph}). At $h=0$, we observe a sequence of topological transitions as $p$ increases: $3\text{Q} \to 5\text{Q} \to 4\text{Q}$. This sequence is a direct consequence of the competition between the 3Q-regime and 4Q-regime interactions. We pinpoint the critical mixing ratios for these transitions as $p_{c1} \approx 0.512$ and $p_{c2} \approx 0.527$, which are in excellent agreement (within 1\% tolerance) with the values reported in Ref.~\cite{Kato2023}.

\subsection{Hidden Topological Transitions}
A central finding of the original study is the existence of topological transitions that are "hidden" from thermodynamic signatures. We verify this behavior by analyzing the specific heat $C$ and the scalar spin chirality $\chi_{sc}$ (Fig.~\ref{fig:hidden_trans}).

As the system transitions from the 3Q to the 5Q state, we observe a clear jump in the monopole charge $N_m$ and a corresponding change in $\chi_{sc}$, signaling a fundamental restructuring of the topological lattice. However, the specific heat $C$ remains smooth across these transition points, exhibiting no singular behavior. This confirms that the transitions are purely topological in nature and do not involve a standard thermodynamic phase transition associated with a symmetry-breaking singularity. The use of the Dual-Search pipeline was critical in resolving these features by avoiding local minima and preserving the adiabatic evolution of the state.
